%==============================================================================
% Section V: Conclusion
%==============================================================================
\section{Conclusion and Future Work}
\label{sec:conclusion}

In this paper, we proposed \textbf{[METHOD NAME]}, a [type of method]
that addresses [the central problem] by [the core idea]. Through
[architectural innovation] and [training innovation], our method achieves
state-of-the-art performance on [N] public benchmarks while being more
efficient than competitive baselines. A detailed ablation study confirms
that each component contributes meaningfully to the overall performance.

% ---- Limitations -----------------------------------------------------------
\paragraph{Limitations.}
Our work has several limitations worth acknowledging.
\textbf{First}, [METHOD NAME] requires [resource / assumption], which may
not hold in [specific real-world setting].
\textbf{Second}, the method exhibits degraded performance on [edge case],
suggesting that further work is needed for [scenario].
\textbf{Third}, our evaluation is limited to [scope]; broader validation
across [other domains] remains an open question.

% ---- Future work -----------------------------------------------------------
\paragraph{Future Work.}
We see several promising directions:
\begin{enumerate}
  \item Extending [METHOD NAME] to [related task / setting], where the
        same insight may yield similar gains.
  \item Combining our approach with [orthogonal technique] to further
        improve [property].
  \item Investigating the theoretical guarantees behind the observed
        [empirical property], which could lead to a deeper understanding
        of [the underlying phenomenon].
\end{enumerate}
We hope this work inspires further research at the intersection of
[area~A] and [area~B].
